\documentclass[10pt,a4paper]{amsart}
\usepackage{amsmath,amssymb,amsthm,graphicx}
% Retained preamble for the Stacks Project chapter snapshot in this directory.
% The chapter's own preamble.tex is a copy of the shared Stacks Project preamble; the block
% below is sliced MECHANICALLY, line for line, from it. Nothing was added, removed or
% redefined; the wrapper class, the xr-hyper/hyperref wiring and the externaldocument list
% are omitted because the assembler supplies the document class.
%
% lines taken from preamble.tex: 3, 5, 9, 11, 12, 156, 157, 158, ... (51 lines total)
%
% For commutative diagrams we use Xy-pic
\usepackage[all]{xy}

% We use 2cell for 2-commutative diagrams.
\xyoption{2cell}
\UseAllTwocells

\newenvironment{reference}{\comment}{\endcomment}
\newenvironment{slogan}{\comment}{\endcomment}
\newenvironment{history}{\comment}{\endcomment}
\theoremstyle{plain}
\newtheorem{theorem}[subsection]{Theorem}
\newtheorem{proposition}[subsection]{Proposition}
\newtheorem{lemma}[subsection]{Lemma}
\theoremstyle{definition}
\newtheorem{definition}[subsection]{Definition}
\newtheorem{example}[subsection]{Example}
\newtheorem{exercise}[subsection]{Exercise}
\newtheorem{situation}[subsection]{Situation}
\theoremstyle{remark}
\newtheorem{remark}[subsection]{Remark}
\newtheorem{remarks}[subsection]{Remarks}
\numberwithin{equation}{subsection}
\def\lim{\mathop{\mathrm{lim}}\nolimits}
\def\colim{\mathop{\mathrm{colim}}\nolimits}
\def\Spec{\mathop{\mathrm{Spec}}}
\def\Hom{\mathop{\mathrm{Hom}}\nolimits}
\def\Ext{\mathop{\mathrm{Ext}}\nolimits}
\def\SheafHom{\mathop{\mathcal{H}\!\mathit{om}}\nolimits}
\def\SheafExt{\mathop{\mathcal{E}\!\mathit{xt}}\nolimits}
\def\Sch{\mathit{Sch}}
\def\Mor{\mathop{\mathrm{Mor}}\nolimits}
\def\Ob{\mathop{\mathrm{Ob}}\nolimits}
\def\Sh{\mathop{\mathit{Sh}}\nolimits}
\def\NL{\mathop{N\!L}\nolimits}
\def\CH{\mathop{\mathrm{CH}}\nolimits}
\def\proetale{{pro\text{-}\acute{e}tale}}
\def\etale{{\acute{e}tale}}
\def\QCoh{\mathit{QCoh}}
\def\Ker{\mathop{\mathrm{Ker}}}
\def\Im{\mathop{\mathrm{Im}}}
\def\Coker{\mathop{\mathrm{Coker}}}
\def\Coim{\mathop{\mathrm{Coim}}}
\DeclareMathSymbol{\boxtimes}{\mathbin}{AMSa}{"02}
\def\QCohstack{\mathcal{QC}\!\mathit{oh}}
\def\Cohstack{\mathcal{C}\!\mathit{oh}}
\def\Spacesstack{\mathcal{S}\!\mathit{paces}}
\def\Quotfunctor{\mathrm{Quot}}
\def\Hilbfunctor{\mathrm{Hilb}}
\def\Curvesstack{\mathcal{C}\!\mathit{urves}}
\def\Polarizedstack{\mathcal{P}\!\mathit{olarized}}
\def\Complexesstack{\mathcal{C}\!\mathit{omplexes}}
\def\Pic{\mathop{\mathrm{Pic}}\nolimits}
\def\Picardstack{\mathcal{P}\!\mathit{ic}}
\def\Picardfunctor{\mathrm{Pic}}
\def\Deformationcategory{\mathcal{D}\!\mathit{ef}}
\usepackage{cleveref}
\crefname{originalreference}{original reference}{original references}
\Crefname{originalreference}{Original reference}{Original references}
\makeatletter
\newlabel{stacks-morphisms-lemma-composition-flat}{{06PX}{source}{}{sourcepointer-1}{}}
\newlabel{stacks-morphisms-lemma-composition-flat@cref}{{[originalreference][0][]06PX}{[1][0][]source}{}{}{}}
\makeatother
\crefname{originalreference}{original reference}{original references}
\Crefname{originalreference}{Original reference}{Original references}
\makeatletter
\newlabel{sites-section-sheafification}{{00W1}{source}{}{sourcepointer-2}{}}
\newlabel{sites-section-sheafification@cref}{{[originalreference][0][]00W1}{[1][0][]source}{}{}{}}
\makeatother
\crefname{originalreference}{original reference}{original references}
\Crefname{originalreference}{Original reference}{Original references}
\makeatletter
\newlabel{homology-lemma-characterize-serre-subcategory}{{02MP}{source}{}{sourcepointer-3}{}}
\newlabel{homology-lemma-characterize-serre-subcategory@cref}{{[originalreference][0][]02MP}{[1][0][]source}{}{}{}}
\makeatother
\crefname{originalreference}{original reference}{original references}
\Crefname{originalreference}{Original reference}{Original references}
\makeatletter
\newlabel{more-morphisms-lemma-etale-dominates-smooth}{{055V}{source}{}{sourcepointer-4}{}}
\newlabel{more-morphisms-lemma-etale-dominates-smooth@cref}{{[originalreference][0][]055V}{[1][0][]source}{}{}{}}
\makeatother
\title{Stability properties of parasitic modules}
\author{The Stacks Project Authors}
\date{stacks-project:stacks-cohomology.tex:local-snapshot-20261012; GNU FDL 1.2 or later}
\begin{document}
\maketitle
\noindent\textit{Editorial note.} All delivered mathematical text is the author's own, copied verbatim from the pinned Stacks Project chapter file stacks-cohomology.tex; no statement or proof step was rewritten, reordered or removed.\par\medskip
\noindent\textit{Allowed standard prerequisites (editorial).}\begin{itemize}
\item Flat composition on stacks; Serre characterization. Morphisms of Stacks, Lemma stacks-morphisms-lemma-composition-flat; Homology, Lemma homology-lemma-characterize-serre-subcategory. Parts (1), (2).
\end{itemize}\medskip
\ifdefined\hypertarget\hypertarget{sourcepointer-1}{}\fi
\noindent\textit{Source reference.} Original reference 06PX: sources/stacks-chow/tags.txt, source lines 8586--8586. Cross-chapter standard prerequisite cited by the author inside the counted span; it is used as a pointer and its own proof is not part of this item.\par\medskip
\ifdefined\hypertarget\hypertarget{sourcepointer-2}{}\fi
\noindent\textit{Source reference.} Original reference 00W1: sources/stacks-chow/tags.txt, source lines 1167--1167. Cross-chapter standard prerequisite cited by the author inside the counted span; it is used as a pointer and its own proof is not part of this item.\par\medskip
\ifdefined\hypertarget\hypertarget{sourcepointer-3}{}\fi
\noindent\textit{Source reference.} Original reference 02MP: sources/stacks-chow/tags.txt, source lines 3452--3452. Cross-chapter standard prerequisite cited by the author inside the counted span; it is used as a pointer and its own proof is not part of this item.\par\medskip
\ifdefined\hypertarget\hypertarget{sourcepointer-4}{}\fi
\noindent\textit{Source reference.} Original reference 055V: sources/stacks-chow/tags.txt, source lines 6688--6688. Cross-chapter standard prerequisite cited by the author inside the counted span; it is used as a pointer and its own proof is not part of this item.\par\medskip
\begin{lemma}
\label{lemma-parasitic}
Let $\mathcal{X}$ be an algebraic stack. Let $\mathcal{F}$
be a presheaf of $\mathcal{O}_\mathcal{X}$-modules.
\begin{enumerate}
\item If $\mathcal{F}$ is parasitic and
$g : \mathcal{Y} \to \mathcal{X}$ is a flat morphism of algebraic stacks,
then $g^*\mathcal{F}$ is parasitic.
\item For $\tau \in \{Zariski, \etale, smooth, syntomic, fppf\}$
we have
\begin{enumerate}
\item the $\tau$ sheafification of a parasitic presheaf of modules is
parasitic, and
\item the full subcategory of
$\textit{Mod}(\mathcal{X}_\tau, \mathcal{O}_\mathcal{X})$
consisting of parasitic modules is a Serre subcategory.
\end{enumerate}
\item Suppose $\mathcal{F}$ is a sheaf for the \'etale topology.
Let $f_i : \mathcal{X}_i \to \mathcal{X}$ be a family of
smooth morphisms of algebraic stacks such that
$|\mathcal{X}| = \bigcup_i |f_i|(|\mathcal{X}_i|)$. If each
$f_i^*\mathcal{F}$ is parasitic then so is $\mathcal{F}$.
\item Suppose $\mathcal{F}$ is a sheaf for the fppf topology.
Let $f_i : \mathcal{X}_i \to \mathcal{X}$ be a family of
flat and locally finitely presented morphisms of algebraic stacks such that
$|\mathcal{X}| = \bigcup_i |f_i|(|\mathcal{X}_i|)$. If each
$f_i^*\mathcal{F}$ is parasitic then so is $\mathcal{F}$.
\end{enumerate}
\end{lemma}

\begin{proof}
To see part (1) let $y$ be an object of $\mathcal{Y}$ which lies
over a scheme $V$ such that the corresponding morphism $y : V \to \mathcal{Y}$
is flat. Then $g(y) : V \to \mathcal{Y} \to \mathcal{X}$ is flat
as a composition of flat morphisms (see
Morphisms of Stacks, Lemma \ref{stacks-morphisms-lemma-composition-flat})
hence $\mathcal{F}(g(y))$ is zero by assumption. Since
$g^*\mathcal{F} = g^{-1}\mathcal{F}(y) = \mathcal{F}(g(y))$ we conclude
$g^*\mathcal{F}$ is parasitic.

\medskip\noindent
To see part (2)(a) note that if $\{x_i \to x\}$ is a $\tau$-covering
of $\mathcal{X}$, then each of the morphisms $x_i \to x$ lies
over a flat morphism of schemes. Hence if $x$ lies over a scheme
$U$ such that $x : U \to \mathcal{X}$ is flat, so do all of the
objects $x_i$. Hence the presheaf $\mathcal{F}^+$ (see
Sites, Section \ref{sites-section-sheafification})
is parasitic if the presheaf $\mathcal{F}$ is
parasitic. This proves (2)(a) as the sheafification of $\mathcal{F}$
is $(\mathcal{F}^+)^+$.

\medskip\noindent
Let $\mathcal{F}$ be a parasitic $\tau$-module. It is immediate from the
definitions that any submodule of $\mathcal{F}$ is parasitic. On the other
hand, if $\mathcal{F}' \subset \mathcal{F}$ is a submodule, then it is
equally clear that the presheaf
$x \mapsto \mathcal{F}(x)/\mathcal{F}'(x)$
is parasitic. Hence the quotient $\mathcal{F}/\mathcal{F}'$ is a parasitic
module by (2)(a). Finally, we have to show that given a short exact sequence
$0 \to \mathcal{F}_1 \to \mathcal{F}_2 \to \mathcal{F}_3 \to 0$
with $\mathcal{F}_1$ and $\mathcal{F}_3$ parasitic, then $\mathcal{F}_2$
is parasitic. This follows immediately on evaluating on $x$ lying
over a scheme flat over $\mathcal{X}$. This proves (2)(b), see
Homology, Lemma \ref{homology-lemma-characterize-serre-subcategory}.

\medskip\noindent
Let $f_i : \mathcal{X}_i \to \mathcal{X}$ be a jointly surjective family of
smooth morphisms of algebraic stacks and assume each $f_i^*\mathcal{F}$
is parasitic. Let $x$ be an object of $\mathcal{X}$ which lies over a
scheme $U$ such that $x : U \to \mathcal{X}$ is flat. Consider a surjective
smooth covering $W_i \to U \times_{x, \mathcal{X}} \mathcal{X}_i$.
Denote $y_i : W_i \to \mathcal{X}_i$ the projection. It follows
that $\{f_i(y_i) \to x\}$ is a covering for the smooth topology
on $\mathcal{X}$. Since a composition of flat morphisms is flat we see that
$f_i^*\mathcal{F}(y_i) = 0$. On the other hand, as we saw in the proof
of (1), we have $f_i^*\mathcal{F}(y_i) = \mathcal{F}(f_i(y_i))$.
Hence we see that for some smooth covering $\{x_i \to x\}_{i \in I}$
in $\mathcal{X}$ we have $\mathcal{F}(x_i) = 0$. This implies
$\mathcal{F}(x) = 0$ because the smooth topology is the same
as the \'etale topology, see
More on Morphisms, Lemma \ref{more-morphisms-lemma-etale-dominates-smooth}.
Namely, $\{x_i \to x\}_{i \in I}$ lies over a smooth covering
$\{U_i \to U\}_{i \in I}$ of schemes. By the lemma just referenced
there exists an \'etale covering $\{V_j \to U\}_{j \in J}$ which
refines $\{U_i \to U\}_{i \in I}$. Denote $x'_j = x|_{V_j}$.
Then $\{x'_j \to x\}$ is an \'etale covering in $\mathcal{X}$
refining $\{x_i \to x\}_{i \in I}$. This means the map
$\mathcal{F}(x) \to \prod_{j \in J} \mathcal{F}(x'_j)$, which is
injective as $\mathcal{F}$ is a sheaf in the \'etale topology,
factors through $\mathcal{F}(x) \to \prod_{i \in I} \mathcal{F}(x_i)$
which is zero. Hence $\mathcal{F}(x) = 0$ as desired.

\medskip\noindent
Proof of (4): omitted. Hint: similar, but simpler, than the proof of (3).
\end{proof}

\end{document}
